
\vspace{-4mm}
\section{Introduction}

Reasoning performance of large language models has advanced rapidly through reinforcement
learning and self-improvement techniques
\citep{shao2024deepseekmath, ouyang2022training, guo2025deepseek}. Among these approaches,
on-policy self-distillation (OPSD) is a promising way to improve reasoning: the student learns
from its own sampled trajectories while a teacher supplies supervision using privileged
information available only during training, such as reference solutions, verified reasoning
traces, or external feedback
\citep{zhao2026self, shenfeld2026self, hubotter2026reinforcement, li2026rethinking,
jang2026stable, yang2026self}. Unlike offline distillation, OPSD lets the teacher respond to
the student's actual intermediate states. Yet this promise has not produced a reliably
effective recipe; in practice, OPSD can be brittle and often demands substantial engineering.

Most existing OPSD methods begin from direct teacher imitation. Given a student-generated
trajectory, the privileged teacher provides token-level predictions, and the student minimizes
a reverse Kullback--Leibler (KL) divergence to those predictions
\citep{zhao2026self, jang2026stable, li2026rethinking, yang2026self}. This formulation fixes
the teacher as the unique optimization target and provides no explicit control over how
aggressively the student should move away from its current or initial policy. Is direct teacher
matching truly the only principled OPSD objective?

Our key observation is that it is not. Vanilla OPSD is precisely the $\beta=1$ member of a
broader KL-regularized policy-optimization family, where $\beta$ weights the KL penalty that
anchors the learned policy to a reference policy. This equivalence turns $\beta$ from an
implicit value fixed at one into an explicit control over the strength of teacher guidance.
It also changes the interpretation of OPSD: rather than a single teacher-imitation objective,
it is one point in a continuum of policy-optimization problems that trade off improvement
toward the privileged teacher against proximity to the reference policy.

This perspective leads to \oursbold, a generalized family of on-policy self-distillation
objectives parameterized by $\beta$. Vanilla OPSD is recovered exactly at $\beta=1$, while
larger values yield more conservative updates that remain closer to the reference policy.
The family therefore exposes a degree of freedom that vanilla OPSD leaves fixed and offers
a principled mechanism for addressing its brittleness: instead of moving abruptly toward
the teacher, the target can progress smoothly from the reference policy toward privileged
teacher guidance over the course of training.

The policy-optimization view also tells us what these targets should be. We derive the
optimal policy of \ours{} in closed form and show that it is a geometric interpolation
between the reference policy and the privileged teacher. Directly optimizing the resulting
RL objective, however, would incur costly rollouts, advantage estimation, and high-variance
gradients; the exact sequence-level optimum also contains an intractable normalizer. We
therefore use policy optimization for the derivation and return to distillation for training.
Each value of $\beta$ defines an RL-derived target along the reference-to-teacher path, which
we realize efficiently through token-level logit interpolation. Scheduling $\beta$ then gives
a smooth curriculum of targets while preserving the computational structure of OPSD. In
short, inexpensive distillation approximates the solution of expensive policy optimization.

Finally, matching the right target still requires assigning credit across the generated
trajectory. Standard token-local OPSD updates are myopic: an early token changes later
prefixes and therefore affects future student--target mismatch. A policy-gradient derivation
of the sequence-level objective yields return-to-go credit assignment, which propagates this
future mismatch back to earlier tokens. Combined with the interpolated target, this gives a
practical algorithm that retains token-level distillation while more faithfully optimizing the
underlying sequence-level objective.

We evaluate \ours{} on mathematical reasoning benchmarks using open-source reasoning
language models ranging from 1.7B to 8B parameters. Across multiple benchmarks, \ours{}
consistently improves optimization stability and reasoning performance over vanilla OPSD.
On Qwen3-1.7B \citep{yang2025qwen3}, it improves avg@12 by up to 9.16 percentage points
over vanilla OPSD across AIME 2024 \citep{aime24}, AIME 2025 \citep{aime25}, and HMMT
2025 \citep{dekoninck2026matharena}. Ablations further isolate the benefits of both the
RL-derived interpolation target and sequence-level return-to-go credit assignment.


Our contributions are summarized as follows: 
\begin{itemize} 
    \item We show that vanilla OPSD is precisely the $\beta=1$ member of a broader
    KL-regularized policy-optimization family, making its implicit regularization choice
    explicit and controllable.
    \item We introduce \oursbold and derive its closed-form optimal policy as a geometric
    interpolation between the reference policy and the privileged teacher.
    \item We turn the RL-derived solution back into an efficient distillation algorithm using
    scheduled logit interpolation and return-to-go credit assignment.
    \item We demonstrate consistent improvements over vanilla OPSD on mathematical
    reasoning benchmarks and validate the contribution of both practical components.
\end{itemize}


\label{sec:intro}
\vspace{-2mm}
